%% Created by Maple 2022.1, Windows 10
%% Source Worksheet: lab_sols08
%% Generated: Wed May 08 10:35:42 BST 2024
\documentclass{article}
\usepackage{amssymb}
\usepackage{graphicx}
\usepackage{hyperref}
\usepackage{listings}
\usepackage{mathtools}
\usepackage{maple}
\usepackage[utf8]{inputenc}
\usepackage[svgnames]{xcolor}
\usepackage{amsmath}
\usepackage{breqn}
\usepackage{textcomp}
\begin{document}
\lstset{basicstyle=\ttfamily,breaklines=true,columns=flexible}
\pagestyle{empty}
\DefineParaStyle{Maple Bullet Item}
\DefineParaStyle{Maple Heading 1}
\DefineParaStyle{Maple Warning}
\DefineParaStyle{Maple Heading 4}
\DefineParaStyle{Maple Heading 2}
\DefineParaStyle{Maple Heading 3}
\DefineParaStyle{Maple Dash Item}
\DefineParaStyle{Maple Error}
\DefineParaStyle{Maple Title}
\DefineParaStyle{Maple Ordered List 1}
\DefineParaStyle{Maple Text Output}
\DefineParaStyle{Maple Ordered List 2}
\DefineParaStyle{Maple Ordered List 3}
\DefineParaStyle{Maple Normal}
\DefineParaStyle{Maple Ordered List 4}
\DefineParaStyle{Maple Ordered List 5}
\DefineCharStyle{Maple 2D Output}
\DefineCharStyle{Maple 2D Input}
\DefineCharStyle{Maple Maple Input}
\DefineCharStyle{Maple 2D Math}
\DefineCharStyle{Maple Hyperlink}
\begin{center}
\begin{Maple Normal}
\underline{\textbf{Integration 2}}
\end{Maple Normal}
\end{center}
\begin{Maple Normal}
\_\_\_\_\_\_\_\_\_\_\_\_\_\_\_\_\_\_\_\_\_\_\_\_\_\_\_\_\_\_\_\_\_\_\_\_\_\_\_\_\_\_\_\_\_\_\_\_\_\_\_\_\_\_\_\_\_\_\_\_\_\_\_\_\_\_\_\_\_\_\_\_\_\_\_\_\_\_\_\_\_
\end{Maple Normal}
\begin{Maple Normal}
\textbf{Exercise 1}
\end{Maple Normal}
\begin{lstlisting}
> p[0] := sqrt(1/2);
p[1] := sqrt(3/2) * x;
p[2] := sqrt(5/8) * (3*x^2 - 1);
p[3] := sqrt(7/8) * (5*x^3 - 3*x);
\end{lstlisting}
\begin{Maple Normal}
We find experimentally that 
${\textcolor{gray}{\int}}_{\!\!-1}^{1}p_{i}(x) p_{j}(x)\textcolor{gray}{d}x$ is 
$0$ unless 
$i =j$, in which case the integral is 
$1$.\

Here are three cases:
\end{Maple Normal}
\begin{lstlisting}
> int(p[1] * p[2],x=-1..1);
\end{lstlisting}
\begin{lstlisting}
> int(p[2] * p[2],x=-1..1);
\end{lstlisting}
\begin{lstlisting}
> int(p[3] * p[2],x=-1..1);
\end{lstlisting}
\begin{Maple Normal}
We can do all 16 cases together and print the results in a table, as follows:
\end{Maple Normal}
\begin{lstlisting}
> Matrix(
 [seq(
   [seq(int(p[i] * p[j],x=-1..1),
        j=0..3)],
  i=0..3)]
);
\end{lstlisting}
\begin{Maple Normal}
We now need to find a function 
$p_{4}(x)=a x^{4}+b x^{2}+c$ that makes the pattern continue: 
\end{Maple Normal}
\begin{lstlisting}
> p[4] := a*x^4 + b*x^2 + c ;
\end{lstlisting}
\begin{Maple Normal}
The following integrals should be 0, 0, 0, 0 and 1, respectively:
\end{Maple Normal}
\begin{lstlisting}
> int(p[0] * p[4],x=-1..1);
int(p[1] * p[4],x=-1..1);
int(p[2] * p[4],x=-1..1);
int(p[3] * p[4],x=-1..1);
int(p[4] * p[4],x=-1..1);
\end{lstlisting}
\begin{Maple Normal}
We can get Maple to find 
$a$, 
$b$ and 
$c$ based on this information:
\end{Maple Normal}
\begin{lstlisting}
> _EnvExplicit := true;
\end{lstlisting}
\begin{lstlisting}
> sol := solve(\{
 int(p[0] * p[4],x=-1..1) = 0,
 int(p[1] * p[4],x=-1..1) = 0,
 int(p[2] * p[4],x=-1..1) = 0,
 int(p[3] * p[4],x=-1..1) = 0,
 int(p[4] * p[4],x=-1..1) = 1,
 a > 0
\},\{a,b,c\});
\end{lstlisting}
\begin{lstlisting}
> p[4] := factor(subs(sol,p[4]));
\end{lstlisting}
\begin{lstlisting}
> q := (n) -> sqrt(n + 1/2)*diff((x^2-1)^n,x$n)/(2^n * n!);
\end{lstlisting}
\begin{lstlisting}
> q(1); q(2); q(3); q(4);
\end{lstlisting}
\begin{lstlisting}
> seq(factor(q(n)/p[n]),n=1..4);
\end{lstlisting}
\begin{Maple Normal}
\_\_\_\_\_\_\_\_\_\_\_\_\_\_\_\_\_\_\_\_\_\_\_\_\_\_\_\_\_\_\_\_\_\_\_\_\_\_\_\_\_\_\_\_\_\_\_\_\_\_\_\_\_\_\_\_\_\_\_\_\_\_\_\_\_\_\_\_\_\_\_\_\_\_\_\_\_\_\_\_\_
\end{Maple Normal}
\begin{Maple Normal}
\textbf{Exercise 2}
\end{Maple Normal}
\begin{Maple Normal}

\end{Maple Normal}
\begin{Maple Normal}
The general picture is that for any positive integers 
$n$ and 
$m$, there are constants 
$a_{0},\mathit{...} ,a_{m}$ such that
\end{Maple Normal}
\begin{center}
$\displaystyle \textcolor{gray}{\int}x^{n} \ln (x)^{m}\textcolor{gray}{d}x =x^{n +1} (a_{0}+a_{1} \ln (x)+\mathit{...} +a_{m} \mathrm{log}(x)^{m})$\end{center}
\begin{Maple Normal}
Here are some examples:
\end{Maple Normal}
\begin{lstlisting}
> factor(int(x^3*ln(x)^5,x));
\end{lstlisting}
\begin{lstlisting}
> factor(int(x^4*ln(x),x));
\end{lstlisting}
\begin{lstlisting}
> factor(int(x^4*ln(x)^3,x));
\end{lstlisting}
\begin{Maple Normal}
\_\_\_\_\_\_\_\_\_\_\_\_\_\_\_\_\_\_\_\_\_\_\_\_\_\_\_\_\_\_\_\_\_\_\_\_\_\_\_\_\_\_\_\_\_\_\_\_\_\_\_\_\_\_\_\_\_\_\_\_\_\_\_\_\_\_\_\_\_\_\_\_\_\_\_\_\_\_\_\_\_
\end{Maple Normal}
\begin{Maple Normal}
\textbf{Exercise 3}
\end{Maple Normal}
\begin{Maple Normal}
These first few integrals involve functions that you may not have seen before.  The 
$\mathrm{erf}()$ function \

is essentially the cumulative probability function for the normal distribution, so is important in\

statistics.  The functions 
$\mathrm{Ei}()$, 
$\mathrm{EllipticF}()$ and 
$\mathrm{hypergeom}()$ are more obscure, but they have\

various applications in more advanced mathematics.  You can learn more about them by entering\

?Ei, ?EllipticF or ?hypergeom in Maple.
\end{Maple Normal}
\begin{Maple Normal}
\textbf{(a)}
\end{Maple Normal}
\begin{lstlisting}
> int(exp(-x^2),x);
\end{lstlisting}
\begin{Maple Normal}
\textbf{(b)}
\end{Maple Normal}
\begin{lstlisting}
> int(1/ln(x),x);
\end{lstlisting}
\begin{Maple Normal}
\textbf{(c)}
\end{Maple Normal}
\begin{lstlisting}
> int(1/(sqrt(1-x^2)*sqrt(1-2*x^2)),x);
\end{lstlisting}
\begin{Maple Normal}
\textbf{(d)}
\end{Maple Normal}
\begin{lstlisting}
> int((x^8+1)^(-1/2),x);
\end{lstlisting}
\begin{Maple Normal}
\textbf{(e)}
\end{Maple Normal}
\begin{lstlisting}
> int(sin(x)*ln(ln(x)),x);
\end{lstlisting}
\begin{Maple Normal}
\textbf{(f)}
\end{Maple Normal}
\begin{lstlisting}
> int(sin(sin(sin(x))),x);
\end{lstlisting}
\begin{Maple Normal}
\textbf{(g)}
\end{Maple Normal}
\begin{lstlisting}
> int(1/sqrt(1+x+x^10),x);
\end{lstlisting}
\begin{Maple Normal}

\end{Maple Normal}
\begin{Maple Normal}
We are looking for a function of 
$x$ that can be defined in three characters, but cannot
\end{Maple Normal}
\begin{Maple Normal}
easily be integrated.  At least one of the characters must be \textcolor[RGB]{255,0,0}{\textbf{x}} (otherwise our function would
\end{Maple Normal}
\begin{Maple Normal}
be constant, so we could integrate it).  That does not leave much space for anything else:
\end{Maple Normal}
\begin{Maple Normal}
we could have a letter or digit, or any of the characters \textcolor[RGB]{255,0,0}{\textbf{+}}, \textcolor[RGB]{255,0,0}{\textbf{-}}, \textcolor[RGB]{255,0,0}{\textbf{*}}, \textcolor[RGB]{255,0,0}{\textbf{/}} or \textcolor[RGB]{255,0,0}{\textbf{^}}.  Most functions 
\end{Maple Normal}
\begin{Maple Normal}
built from these ingredients are easy to integrate:
\end{Maple Normal}
\begin{lstlisting}
> int(x+3,x);
\end{lstlisting}
\begin{lstlisting}
> int(x*x,x);
\end{lstlisting}
\begin{lstlisting}
> int(2^x,x);
\end{lstlisting}
\begin{Maple Normal}
There is just one exception:
\end{Maple Normal}
\begin{lstlisting}
> int(x^x,x);
\end{lstlisting}
\begin{Maple Normal}
There is actually an answer that is even shorter:
\end{Maple Normal}
\begin{lstlisting}
> int(x!,x);
\end{lstlisting}
\begin{Maple Normal}
This is a little bit dubious because 
$x !$ is usually only defined when 
$x$ is an integer, and that would make\

the integral meaningless.  However, Maple does have a way to interpret 
$x !$ even when 
$x$ is not an \

integer (eg 
$ 1.5!= 1.329340388$), which you can learn about by entering ?factorial and ?GAMMA.
\end{Maple Normal}
\begin{lstlisting}
> 
\end{lstlisting}
\begin{Maple Normal}
\_\_\_\_\_\_\_\_\_\_\_\_\_\_\_\_\_\_\_\_\_\_\_\_\_\_\_\_\_\_\_\_\_\_\_\_\_\_\_\_\_\_\_\_\_\_\_\_\_\_\_\_\_\_\_\_\_\_\_\_\_\_\_\_\_\_\_\_\_\_\_\_\_\_\_\_\_\_\_\_\_
\end{Maple Normal}
\begin{Maple Normal}
\textbf{Exercise 4}
\end{Maple Normal}
\begin{lstlisting}
> y := sin(x) + sin(3*x)/3 + sin(5*x)/5 + sin(7*x)/7;
\end{lstlisting}
\begin{lstlisting}
> z := int(y,x);
\end{lstlisting}
\begin{Maple Normal}
The graph of 
$y$ spends half the time wiggling near 
$ 0.8$ and half the time wiggling near 
$- 0.8$, jumping\

rapidly between these levels at 
$x =0$, 
$x =\mathrm{Pi}$, 
$x =2 \mathrm{Pi}$ and all other multiples of 
$\mathrm{Pi}$.  The graph of\


$z$ is a sawtooth, climbing from about 
$- 1.2$ at 
$x =0$ in a roughly straight line to about 
$ 1.2$ at\


$x =\mathrm{Pi}$, then dropping in a roughly straight line to about 
$- 1.2$ again at 
$x =2 \mathrm{Pi}$, and so on.  This\

makes sense because 
$y =\frac{\mathit{dz}}{\mathit{dx}}$, which is the slope of the graph of 
$z$.  When 
$y$ is near 
$ 0.8$, the graph\

of 
$z$ slopes upwards at a roughly constant angle; when 
$y$ is near 
$- 0.8$, the graph of 
$z$ slopes down\

at a nearly constant angle. 
\end{Maple Normal}
\begin{lstlisting}
> plot([y,z],x=-2*Pi..2*Pi);
\end{lstlisting}
\begin{lstlisting}
> unassign('y','z');
\end{lstlisting}
\begin{Maple Normal}
\_\_\_\_\_\_\_\_\_\_\_\_\_\_\_\_\_\_\_\_\_\_\_\_\_\_\_\_\_\_\_\_\_\_\_\_\_\_\_\_\_\_\_\_\_\_\_\_\_\_\_\_\_\_\_\_\_\_\_\_\_\_\_\_\_\_\_\_\_\_\_\_\_\_\_\_\_\_\_\_\_
\end{Maple Normal}
\begin{Maple Normal}
\textbf{Exercise 5}
\end{Maple Normal}
\begin{lstlisting}
> y := x*sin(20*x)*exp(-x);
\end{lstlisting}
\begin{lstlisting}
> z := 20*int(y,x);
\end{lstlisting}
\begin{Maple Normal}
In this plot, 
$y$ and 
$z$ look very similar:
\end{Maple Normal}
\begin{lstlisting}
> plot([y,z],x=0..10);
\end{lstlisting}
\begin{Maple Normal}
When we look more closely, we see that size of the oscillations in the two graphs are nearly the\

same, as is the frequency of the oscillations, but the oscillations in 
$y$ lag behind those in 
$z$.
\end{Maple Normal}
\begin{lstlisting}
> plot([y,z],x=0..3);
\end{lstlisting}
\begin{Maple Normal}
By focusing still further, we see that the extent of the lag is about a quarter of a complete cycle.
\end{Maple Normal}
\begin{lstlisting}
> plot([y,z],x=1.255..1.57);
\end{lstlisting}
\begin{Maple Normal}
All this makes sense if we look at the formulae:
\end{Maple Normal}
\begin{lstlisting}
> y;
\end{lstlisting}
\begin{lstlisting}
> z;
\end{lstlisting}
\begin{Maple Normal}
If we expand out 
$z$, the first term is 
$-(\frac{400}{401}) x \,{\mathrm e}^{-x} \cos (20 x)$, which is nearly 
$-x \,{\mathrm e}^{-x} \cos (20 x)$.  \

The remaining terms are much smaller.  The function 
$-\cos (20 x)$ is the same as 
$\sin (20 x)$ shifted by a \

quarter cycle (because 
$\sin (\theta -\frac{\mathrm{Pi}}{2})=-\cos (\theta)$).  Thus, everything fits together.\

\


\end{Maple Normal}
\begin{Maple Normal}
In the first step mentioned above, you might want to enter \textcolor[RGB]{255,0,0}{\textbf{expand(z)}} rather than doing it by hand.\

This does not work very well, as Maple insists on converting 
$\cos (20 x)$ into a polynomial in 
$\cos (x)$,\

giving a very messy result.  However, the following (rather obscure) syntax does the trick:
\end{Maple Normal}
\begin{lstlisting}
> evalf(frontend('expand',[z]));
\end{lstlisting}
\begin{Maple Normal}

\end{Maple Normal}
\begin{Maple Normal}
We now consider 
$y^{2}+z^{2}$:
\end{Maple Normal}
\begin{lstlisting}
> plot(y^2+z^2,x=0..6);
\end{lstlisting}
\begin{lstlisting}
> simplify(y^2+z^2);
\end{lstlisting}
\begin{Maple Normal}
\

The first term is 
$(\frac{25921282001}{25856961601}) x^{2} {\mathrm e}^{-2 x}$.  As 
$25921282001$ is close to 
$25856961601$, \

this is approximately 
$x^{2} {\mathrm e}^{-2 x}$.  The other terms are much smaller, so we guess that 
$y^{2}+z^{2}$\

should be close to 
$x^{2} {\mathrm e}^{-2 x}$.  We can check this graphically as follows:
\end{Maple Normal}
\begin{lstlisting}
> plot([y^2+z^2,x^2*exp(-2*x)],x=0..6);
\end{lstlisting}
\begin{lstlisting}
> unassign('y','z');
\end{lstlisting}
\begin{Maple Normal}
\_\_\_\_\_\_\_\_\_\_\_\_\_\_\_\_\_\_\_\_\_\_\_\_\_\_\_\_\_\_\_\_\_\_\_\_\_\_\_\_\_\_\_\_\_\_\_\_\_\_\_\_\_\_\_\_\_\_\_\_\_\_\_\_\_\_\_\_\_\_\_\_\_\_\_\_\_\_\_\_\_
\end{Maple Normal}
\begin{Maple Normal}
\textbf{Exercise 6}
\end{Maple Normal}
\begin{Maple Normal}
Here is a solution in a single step:
\end{Maple Normal}
\begin{lstlisting}
> solve(
  \{seq(int(x^k*(a*x^2+b*x+c)*exp(-x),x=0..infinity)=k,k=1..3)\},
  \{a,b,c\}
);
\end{lstlisting}
\begin{Maple Normal}
Here it is again, more slowly:
\end{Maple Normal}
\begin{lstlisting}
> y := (a*x^2+b*x+c)*exp(-x);
\end{lstlisting}
\begin{lstlisting}
> int1 := int(x*y,x=0..infinity);
\end{lstlisting}
\begin{lstlisting}
> int2 := int(x^2*y,x=0..infinity);
\end{lstlisting}
\begin{lstlisting}
> int3 := int(x^3*y,x=0..infinity);
\end{lstlisting}
\begin{lstlisting}
> solve(\{int1 = 1, int2 = 2, int3 = 3\},\{a,b,c\});
\end{lstlisting}
\begin{lstlisting}
> unassign('y');
\end{lstlisting}
\begin{Maple Normal}
\_\_\_\_\_\_\_\_\_\_\_\_\_\_\_\_\_\_\_\_\_\_\_\_\_\_\_\_\_\_\_\_\_\_\_\_\_\_\_\_\_\_\_\_\_\_\_\_\_\_\_\_\_\_\_\_\_\_\_\_\_\_\_\_\_\_\_\_\_\_\_\_\_\_\_\_\_\_\_\_\_
\end{Maple Normal}
\begin{Maple Normal}
\textbf{Exercise 7}
\end{Maple Normal}
\begin{lstlisting}
> fsolve(x*(ln(ln(x)) + ln(x) - 1) = 541,\{x=100\});
\end{lstlisting}
\begin{Maple Normal}
\_\_\_\_\_\_\_\_\_\_\_\_\_\_\_\_\_\_\_\_\_\_\_\_\_\_\_\_\_\_\_\_\_\_\_\_\_\_\_\_\_\_\_\_\_\_\_\_\_\_\_\_\_\_\_\_\_\_\_\_\_\_\_\_\_\_\_\_\_\_\_\_\_\_\_\_\_\_\_\_\_
\end{Maple Normal}
\begin{Maple Normal}
\textbf{Exercise 8}
\end{Maple Normal}
\begin{lstlisting}
> seq(a+b^n,n=2..10);
\end{lstlisting}
\begin{lstlisting}
> 
\end{lstlisting}
\begin{Maple Normal}
\_\_\_\_\_\_\_\_\_\_\_\_\_\_\_\_\_\_\_\_\_\_\_\_\_\_\_\_\_\_\_\_\_\_\_\_\_\_\_\_\_\_\_\_\_\_\_\_\_\_\_\_\_\_\_\_\_\_\_\_\_\_\_\_\_\_\_\_\_\_\_\_\_\_\_\_\_\_\_\_\_
\end{Maple Normal}
\begin{Maple Normal}
\textbf{Exercise 9}
\end{Maple Normal}
\begin{lstlisting}
> with(plots):
\end{lstlisting}
\begin{lstlisting}
> display(
 implicitplot(x^4+y^4=1,x=-2..2,y=-2..2),
 plot([cos(t)/(2+cos(4*t)),sin(t)/(2+cos(4*t)),t=0..2*Pi])
);
\end{lstlisting}
\begin{lstlisting}
> 
\end{lstlisting}
\begin{lstlisting}
> 
\end{lstlisting}
\begin{Maple Normal}
\_\_\_\_\_\_\_\_\_\_\_\_\_\_\_\_\_\_\_\_\_\_\_\_\_\_\_\_\_\_\_\_\_\_\_\_\_\_\_\_\_\_\_\_\_\_\_\_\_\_\_\_\_\_\_\_\_\_\_\_\_\_\_\_\_\_\_\_\_\_\_\_\_\_\_\_\_\_\_\_\_
\end{Maple Normal}
\begin{Maple Normal}
\textbf{Exercise 10}
\end{Maple Normal}
\begin{lstlisting}
> evalf[100](Pi^76/exp(87));
\end{lstlisting}
\begin{lstlisting}
> 
\end{lstlisting}
\begin{Maple Normal}
\_\_\_\_\_\_\_\_\_\_\_\_\_\_\_\_\_\_\_\_\_\_\_\_\_\_\_\_\_\_\_\_\_\_\_\_\_\_\_\_\_\_\_\_\_\_\_\_\_\_\_\_\_\_\_\_\_\_\_\_\_\_\_\_\_\_\_\_\_\_\_\_\_\_\_\_\_\_\_\_\_
\end{Maple Normal}
\begin{Maple Normal}
\textbf{Exercise 11}
\end{Maple Normal}
\begin{lstlisting}
> simplify(sin(2*x) * tan(2*x) * diff(log(tan(x)),x,x));
\end{lstlisting}
\begin{lstlisting}
> 
\end{lstlisting}
\begin{lstlisting}
> 
\end{lstlisting}
\end{document}
